\documentclass[10pt]{article}
\usepackage[margin=1in]{geometry}
\usepackage[T1]{fontenc}
\usepackage{lmodern,microtype,amsmath,amssymb,amsthm}
\theoremstyle{definition}
\newtheorem{assumption}{Assumption}
\begin{document}
\section*{Adaptive Gradient Descent with Lipschitz Hessian}
We consider minimizing a possibly nonconvex, twice continuously differentiable function $f:E\to\mathbb{R}$ on a finite-dimensional real inner-product space $E$. The gradient and Hessian are Lipschitz with constants $L>0$ and $\rho\geq0$, respectively. A given starting point $x^0\in E$ satisfies $f(x^0)-\inf_{x\in E}f(x)\leq\Delta$, where $\Delta\geq0$. The oracle returns $\nabla f(x)$ at $x$. Adaptive gradient descent must return a queried point $\bar{x}$ satisfying $\|\nabla f(\bar{x})\|\leq\varepsilon$, for $\varepsilon>0$, almost surely under $P$, the law of the algorithm's randomness; with no queries it returns $x^0$.

The optimization problem is
\[
\inf_{x\in E} f(x).
\]
The norm $\|\cdot\|$ is induced by the inner product on $E$, and $\|\cdot\|_{\mathrm{op}}$ denotes the induced operator norm. The Hessian $\nabla^2 f(x)$ acts on $E$. The algorithm receives $E$, $x^0$, $L$, $\rho$, $\Delta$, and $\varepsilon$.

Our assumptions:
\begin{assumption}[Smoothness and Hessian regularity]
The function $f$ is twice continuously Fr\'echet differentiable on $E$. For every $x,y\in E$,
\[
\begin{aligned}
\|\nabla f(x)-\nabla f(y)\|&\leq L\|x-y\|,\\
\|\nabla^2 f(x)-\nabla^2 f(y)\|_{\mathrm{op}}
&\leq \rho\|x-y\|.
\end{aligned}
\]
\end{assumption}

\begin{assumption}[Initial objective gap]
The given starting point satisfies
\[
f(x^0)-f(x)\leq\Delta,
\qquad x\in E.
\]
\end{assumption}

Our formulation is algorithm-specific: we analyze adaptive gradient descent with strictly positive step sizes and an output selected by the smallest observed gradient norm. Queries may be made at any point of $E$. At iteration $k$, the method queries $x^k$, receives $g^k=\nabla f(x^k)$, and updates
\[
x^{k+1}=x^k-\eta_k g^k,
\qquad \eta_k>0.
\]
The step size may depend on the algorithm's current information, its randomness, and the current reply. Its step-selection rule must return a strictly positive value for every possible reply.

To specify the output exactly, let $b^k\in E$ be the stored point and $s^k\in\mathbb{R}$ its stored squared gradient norm. Initialize $b^0=x^0$; the initial value $s^0$ is unrestricted. After each reply, set
\[
\begin{aligned}
b^{k+1}&=
\begin{cases}
x^k, & k=0\ \text{or}\ \|g^k\|^2\leq s^k,\\
b^k, & \text{otherwise},
\end{cases}\\
s^{k+1}&=
\begin{cases}
\|g^k\|^2, & k=0\ \text{or}\ \|g^k\|^2\leq s^k,\\
s^k, & \text{otherwise}.
\end{cases}
\end{aligned}
\]
After an allotted number $K\geq0$ of oracle calls, the output is $\bar{x}=b^K$. For $K\geq1$, this is the most recently queried point attaining the smallest gradient norm among $x^0,\ldots,x^{K-1}$. The required accuracy is
\[
\|\nabla f(\bar{x})\|\leq\varepsilon
\quad P\text{-a.s.}
\]
\end{document}
